% Typeset using lualatex

\documentclass[a4paper,10pt]{article}
\usepackage[utf8]{luainputenc}
\usepackage[pdfborder={0 0 0}]{hyperref}

\usepackage[czech]{babel}

\usepackage{fouriernc}
\usepackage{tgpagella}
\usepackage{xcolor}
\usepackage[activate=true]{microtype}
\usepackage{amssymb}
\usepackage{fontspec}

\usepackage{calc}
\usepackage{multicol}
\usepackage{array}
\usepackage{enumitem}
\usepackage{tikz}
\usepackage{siunitx}
\usepackage{colortbl}
\usepackage[a4paper,margin=2cm,top=1cm]{geometry}

\setmainfont[Numbers=OldStyle]{TeX Gyre Pagella}

\definecolor{plhome}{HTML}{CCEEFF}
\definecolor{silver}{HTML}{DDDDDD}

\sisetup{
    output-decimal-marker={,}
}
\pagestyle{empty}

\parskip=1ex
\parindent=0pt

\setlist[enumerate,1]{start=0}

\newcommand\plpage{
    \newpage
    \begin{tikzpicture}[remember picture,overlay]
    \node[opacity=0.1,below left,xshift=1cm,yshift=1cm] at (current page.north east) {\includegraphics[width=12cm]{../../images/pylady-pink}};
    \end{tikzpicture}
}

\newcommand\answerspace{\\\rule[0cm]{0pt}{1cm}}

\newcommand\True{\texttt{True}}
\newcommand\False{\texttt{False}}

\begin{document}

\plpage

\section*{Domácí úkoly 2}

\begin{enumerate}
\item Jakou používáš verzi Pythonu?

\item Jak se jmenuje druh chyby, která nastane, když...
    \begin{itemize}
        \item Zapomeneš dát uvozovky kolem řetězce? — \texttt{SyntaxError}
        \item Zkusíš sečíst číslo a řetězec? —
        \item Dělíš nulou? —
        \item Použiješ proměnnou, do které není nic přiřazené? —
        \item Stiskneš Ctrl+C, když se program ptá na vstup (pomocí \texttt{input})? —
        \item Odsadíš příkaz bez předchozího \texttt{if:}? —
        \item Po \texttt{if:} odsadíš jeden příkaz o čtyři mezery, a druhý jen o dvě? —
        \item Neuzavřeš závorku? –
        \item Neuzavřeš uvozovky u řetězce? –
        \item Zkusíš použít vykřičník (\texttt{!}) jako operátor? —
        \item Napíšeš v příkazu \texttt{print(1, 2, 3)} čárku navíc? —
    \end{itemize}

\item \emph{\small Podívej se na odpovědi na předchozí otázku, ale Python zkus použít jen na ověření.}
    \\Jaká chyba nastane, když zkusíš odečíst číslo od řetězce?

\item \emph{\small Podívej se na odpovědi na otázku 1, ale Python použij jen na ověření.}
    \\Jaká chyba nastane, když zkusíš použít proměnnou předtím, než do ní něco přiřadíš?

\item \emph{\small Podívej se na odpovědi na otázku 1, ale Python použij jen na ověření.}
    \\Jaká chyba nastane, když zkusíš podělit řetězec řetězcem?

\item Lomítko (\texttt{/}) je operátor, ale nedá se použít na řetězce. Vykřičník (\texttt{!}) v Pythonu není operátor.
    \\Jak se liší „jejich“ chyby? \emph{\small Rozdíly jsou aspoň dva.}
    \\Je zavináč (\texttt{@}) v Pythonu operátor?  % odpověď na tohle se v Pythonu 3.5 změní. To bude sranda!

\end{enumerate}

\hrule

\begin{enumerate}[resume]
\item Ne všechno se dá použít jako jméno proměnné.
    \\Fungují pro proměnné následující jména? Pokud ne, proč asi?
    {
        \newcommand\varname[2][]{\varnameend[#1]{#2}\\[0.5cm]}
        \newcommand\varnameend[2][]{\texttt{#2} #1}
        \begin{multicols}{4}
        \varname{x}
        \varname{tlacitko4}
        \varname{34}
        \varname{3e4}
        \varname{krůta}
        \varname{\$i}
        \varname{druha-odmocnina}
        \varname{pet.svestek}
        \varname{kratsiStrana}
        \varname{POCET\_BODU}
        \varname[(podtržítko)]{\_}
        \varname[(pí)]{\textrm{π}}
        \varname{True}
        \varname{\_cache}
        \varname{\_\_name\_\_}
        \varnameend{while}
        \end{multicols}
    }
    \hspace{0.5cm}
\item Jaké nové druhy chyb jsi v předchozím úkolu viděla?

\end{enumerate}

\hrule

\begin{enumerate}[resume]

\item \emph{\small Zkus se nedívat na programy ze srazu.}
    \\Napiš program, který spočítá povrch a obsah krychle o straně \SI{2852}{cm}.
    \\Abys nemusela tolik hledat v učebnici (vlastně Wikipedii): povrch $S=6a^2$, obsah $V=a^3$
    \\Řešení, pro kontrolu: $S=\SI{48803424}{cm^2}$, $V=\SI{23197894208}{cm^3}$

\item Doplň program z minulého úkolu, aby počítal i povrch a obsah koule o poloměru \SI{2852}{cm}.
    \\Vzorečky: povrch $S=4\pi r^2$, obsah $V=\frac34\pi r^3$, $\pi \approx 3.1415926$
    \\Řešení, pro kontrolu: $S\approx\SI{102213650,46}{cm^2}$, $V\approx\SI{54658749584,57}{cm^3}$

\item Použila jsi operátor pro umocnění (\texttt{**})? Jestli ne, naprav to.

\item Změň program tak, aby stranu/poloměr mohl uživatel zadat.

\end{enumerate}

\plpage

\begin{enumerate}[resume]

\item Které Pythoní operátory dokážeš z hlavy vyjmenovat?

\item Zkusila jsi porovnávat řetězce?
    \\Doplň tuhle tabulku tužkou; pak ověř odpovědi pomocí Pythonu:

    {
        \newcommand\rowend{\rule{0pt}{0.5cm}\\ \hline}
        \begin{tabular}{c|c|c|c|c|c}
        \arrayrulecolor{silver}
        \verb+a+ & \verb+b+ & \verb+a == b+ & \verb+a < b+ & \verb+a >= b+ \\
        \arrayrulecolor{black}\hline\arrayrulecolor{silver}
        \verb+1+ & \verb+1+ & \True & \False & \True \rowend
        \verb+1+ & \verb+2+ & \False & \True & \False \rowend
        \verb+'abc'+ & \verb+'abc'+ & \True & \False & \True \rowend
        \verb+'aaa'+ & \verb+'abc'+ & & & \rowend
        \verb+'abc'+ & \verb+'Abc'+ & & & \rowend
        \verb+'abC'+ & \verb+'abc'+ & & & \rowend
        \verb+'abc'+ & \verb+'abcde'+ & & & \rowend
        \verb+'abc'+ & \verb+'ábč'+ & & & \rowend
        \verb+'abc'+ & \verb+10+ & & & \rowend
        \end{tabular}
    }

\item Jaká je hodnota proměnné po provedení příkazu: \hspace{1cm}\verb+promenna = 1 < 2+ \hspace{0.2cm}?
    \\Co když do proměnné přiřadíš \verb+2 < 2+ ?
    \\\verb+1 < 2 < 3+ ?
    \\\verb+1 < 3 < 2+ ?
    \\\verb+1 < 3 < 3+ ?
    \\\verb+'abc' < 'ABC' < 'def' < 'zajíc'+ ?
    \\Pokud je hodnota proměnné \False, dají se některé z porovnávaných hodnot vyměnit, aby byla \True?

\end{enumerate}

\hrule

\begin{enumerate}[resume]

\item Na srazu jsme měli program, který píše různé nesmysly podle uživatelem zadaného věku.
    \\Zkus napsat program, který píše hlášky podle zadané
    rychlosti chůze, váhy ulovené ryby, počtu tykadel, teploty vody
    nebo třeba vzdálenosti od rovníku.
    \\Program pošli na soukromý e-mail organizátora (ne do diskusní skupiny).
    Pošli ho jako přílohu, nekopíruj ho do textu e-mailu.

\item Napiš program, který po zadání správného hesla vypíše nějakou tajnou informaci.
    \\\emph{\small Vhodné tajemství je třeba: V pátek jsem viděla černého havrana.}

\item Projdi si \url{http://python.cz/pyladies/s002-hello-world/and-or.html}

\item Doplň tuhle tabulku:

    {
        \newcommand\rowend{\rule{0pt}{0.5cm}\\ \hline}
        \begin{tabular}{c|c|c|c|c}
        \arrayrulecolor{silver}
        \verb+a+ & \verb+b+ & \verb+a and b+ & \verb+a or b+ & \verb+not a+ \\
        \arrayrulecolor{black}\hline\arrayrulecolor{silver}
        \True & \True &  &  &  \rowend
        \False & \True &  &  &  \rowend
        \True & \False &  &  &  \rowend
        \False & \False &  &  &  \rowend
        \end{tabular}
    }

    \emph{\small Tohle je takzvaná \emph{pravdivostní tabulka}.
          Obsahuje všechny kombinace booleovských hodnot, které \texttt{a} a \texttt{b}
          můžou mít.}

\item Zkus program „šťastná/bohatá“ přepsat pomocí zanořených \texttt{if}ů. Která verze ti připadá čitelnější?

\item Zkus přepsat Kámen, Nůžky, Papír pomocí \texttt{and} a \texttt{or}.
    \\Dokážeš docílit toho, aby se každý z řetězců \texttt{'Plichta.'},
    \texttt{'Počítač vyhrál.'} a \texttt{'Vyhrála jsi!'} objevil v programu jen jednou,
    aniž bys tyhle řetězce musela přiřazovat do proměnných?
    \\Pokud ano, gratuluji!

\end{enumerate}


\end{document}

